\documentclass[11pt]{article}

\usepackage[margin=0.6in]{geometry}
\usepackage{amsmath}
\usepackage{amssymb}
\usepackage{array}
\usepackage{enumitem}
\usepackage{textcomp}
\usepackage{parskip}

\setlist[enumerate,2]{label=\Alph*., leftmargin=2.2em, topsep=2pt, itemsep=1pt}
\setlist[enumerate,1]{leftmargin=2.2em, itemsep=6pt}

\pagestyle{plain}

% Question worth 2 points each (15 questions x 2p = 30p per subject)
\newcommand{\pts}{\textbf{(2p)}~}

% Student name field, shown top-left of each subject
\newcommand{\nameheader}{%
  \noindent\textbf{Name:}~\underline{\hspace{8cm}}\par\vspace{4pt}
}

% Compact centered title
\newcommand{\examtitle}{%
  \begin{center}
  {\large\bfseries Mobile and Embedded Computing; Final Exam (30p)}
  \end{center}
  \vspace{-4pt}
}

\newcommand{\answergrid}{%
  \begin{center}
  \renewcommand{\arraystretch}{1.6}
  \begin{tabular}{|c|*{15}{c|}}
  \hline
  Answer & 1 & 2 & 3 & 4 & 5 & 6 & 7 & 8 & 9 & 10 & 11 & 12 & 13 & 14 & 15 \\
  \hline
  A &&&&&&&&&&&&&&& \\
  \hline
  B &&&&&&&&&&&&&&& \\
  \hline
  C &&&&&&&&&&&&&&& \\
  \hline
  D &&&&&&&&&&&&&&& \\
  \hline
  E &&&&&&&&&&&&&&& \\
  \hline
  \end{tabular}
  \end{center}
}

% Filled answer key for Variant 1
\newcommand{\answergridone}{%
  \begin{center}
  \renewcommand{\arraystretch}{1.6}
  \begin{tabular}{|c|*{15}{c|}}
  \hline
  Answer & 1 & 2 & 3 & 4 & 5 & 6 & 7 & 8 & 9 & 10 & 11 & 12 & 13 & 14 & 15 \\
  \hline
  A & X & X & X & X & X & X & X &   & X & X & X & X & X & X & X \\
  \hline
  B & X & X & X &   &   & X & X & X & X & X & X & X & X & X & X \\
  \hline
  C & X & X & X & X & X & X & X & X &   & X & X & X & X & X & X \\
  \hline
  D &   &   &   & X & X & X & X & X & X & X &   &   &   &   &   \\
  \hline
  E &   & X & X & X & X &   &   & X & X &   & X & X & X &   & X \\
  \hline
  \end{tabular}
  \end{center}
}

% Filled answer key for Variant 2
\newcommand{\answergridtwo}{%
  \begin{center}
  \renewcommand{\arraystretch}{1.6}
  \begin{tabular}{|c|*{15}{c|}}
  \hline
  Answer & 1 & 2 & 3 & 4 & 5 & 6 & 7 & 8 & 9 & 10 & 11 & 12 & 13 & 14 & 15 \\
  \hline
  A & X & X & X & X & X & X & X & X & X & X & X & X & X & X & X \\
  \hline
  B & X & X & X & X & X & X & X & X & X & X & X & X & X & X & X \\
  \hline
  C & X & X & X & X & X & X & X &   &   & X & X & X & X & X & X \\
  \hline
  D &   &   &   &   &   &   & X & X & X &   &   &   &   & X & X \\
  \hline
  E & X & X & X & X &   & X &   & X & X & X & X & X & X &   &   \\
  \hline
  \end{tabular}
  \end{center}
}

% Exam instructions + how to fill in the answer grid
\newcommand{\instructions}{%
  \small
  \noindent Each question has \textbf{zero or more} correct answers and is worth \textbf{2p}. A question is awarded points only if \emph{all} correct options are selected and \emph{no} incorrect options are selected (no partial credit). A non-programmable calculator is allowed for the calculation questions.

  \smallskip
  \noindent\textbf{How to fill in your answers.} All answers go in the grid above. The \textbf{columns (1--15)} correspond to the questions and the \textbf{rows (A--E)} correspond to the options. For each question, put an \textbf{X} in the box where that question's column meets the row of \emph{every} option you consider correct, and leave the boxes of the incorrect options empty (for example, if you think A and C are correct for question 4, mark an X in column 4, rows A and C only). Only the grid is graded. Any marks made next to the questions themselves are ignored. To revise an answer, completely fill the mistaken box and clearly X the intended one.
  \normalsize\par\smallskip
}

\begin{document}

  \nameheader
  \examtitle

%==================================================================
  \subsection*{Variant 1}
%==================================================================

  \answergridone

  \instructions

  \begin{enumerate}

    \item \pts A frame budget is $1000\,\text{ms} \div \text{target FPS}$. Assume the UI thread does all rendering work. Select all that are TRUE:
    \begin{enumerate}
      \item At 60 FPS the per-frame budget is about 16.67 ms ($1000 \div 60$).
      \item At 120 FPS the per-frame budget is about 8.33 ms ($1000 \div 120$).
      \item A synchronous task that runs for 100 ms on the UI thread blocks roughly 6 frames at 60 FPS ($100 \div 16.67 \approx 6$).
      \item At 30 FPS the per-frame budget is about 8.33 ms.
      \item A task that takes exactly 16.67 ms leaves about 50\% of the 60 FPS frame budget free for other work.
    \end{enumerate}

    \item \pts A client retries with exponential backoff: the base delay is \textbf{1 s} and it \textbf{doubles} every attempt (1, 2, 4, 8, \dots), with the delay computed for the \emph{n}-th retry as $2^{\,n-1}$ seconds. Select all that are TRUE:
    \begin{enumerate}
      \item The uncapped backoff delay before the 5th retry is 16 s.
      \item With a 30 s maximum-delay cap, the 6th retry's delay is capped at 30 s instead of its uncapped value of 32 s.
      \item Using \textbf{full jitter} on a computed backoff delay of 8 s, the actual wait is a random value in the range [0 s, 8 s].
      \item Using \textbf{equal jitter} on a computed backoff delay of 8 s, the actual wait is a random value in the range [8 s, 16 s].
      \item Summing the uncapped delays of the first four retries (1 + 2 + 4 + 8) gives 15 s of total waiting.
    \end{enumerate}

    \item \pts Select all that are TRUE:
    \begin{enumerate}
      \item A PWA must be served over HTTPS.
      \item A PWA relies on a service worker that checks for updates in the background and applies them on the next launch.
      \item Native apps generally offer the best performance and use native components.
      \item Cross-platform apps never need to go through a store review process for any kind of update.
      \item Flutter, as a cross-platform framework, also supports Desktop targets.
    \end{enumerate}

    \item \pts Select all that are TRUE:
    \begin{enumerate}
      \item Compiled languages tend to be roughly 100 times faster than interpreted languages.
      \item Interpreted languages catch more errors at compile time, preventing them from reaching production.
      \item Compiled code is generally more difficult to reverse engineer or modify.
      \item Interpreted languages offer a shorter development feedback loop, well-suited to prototyping.
      \item With a compiled language you must ``rebuild'' the program every time you make a change.
    \end{enumerate}

    \item \pts Select all that are TRUE:
    \begin{enumerate}
      \item Dart is null safe.
      \item Dart is dynamically typed and has no static type system.
      \item Dart is garbage collected.
      \item Dart supports hot reload.
      \item Dart has built-in support for asynchronous programming.
    \end{enumerate}

    \item \pts Select all that are TRUE:
    \begin{enumerate}
      \item The Dart GC uses a two-generation approach: a young generation and an old generation.
      \item A remembered set tracks references from old-generation objects to young-generation objects.
      \item A write barrier is code that runs whenever the application modifies an object reference.
      \item Objects that survive multiple young-generation collections are promoted to the old generation.
      \item The remembered set primarily tracks references that point from the young generation into the old generation.
    \end{enumerate}

    \item \pts A write barrier adds information to the remembered set, and the Dart VM has particular execution modes. Select all that are TRUE:
    \begin{enumerate}
      \item The barrier triggers a remembered-set update when the \emph{target} object being written to is in the old generation.
      \item The barrier triggers a remembered-set update when the \emph{new reference} points to a young-generation object.
      \item During development the Dart VM operates in JIT mode.
      \item In production the Dart VM runs in AOT mode (compiled ahead-of-time to native machine code).
      \item AOT compilation increases VM overhead and causes higher memory consumption than JIT.
    \end{enumerate}

    \item \pts Select all that are TRUE:
    \begin{enumerate}
      \item Isolates share a single memory heap and therefore require locks/mutexes for synchronization.
      \item Each isolate has its own memory heap and its own event loop.
      \item Isolates communicate through message passing via ports.
      \item Because isolates do not share mutable memory, race conditions are effectively impossible between them.
      \item Messages containing primitive types are copied when sent between isolates.
    \end{enumerate}

    \item \pts Select all that are TRUE:
    \begin{enumerate}
      \item The asynchronous model is the better fit for I/O-bound operations such as network calls.
      \item Threads/isolates are the better fit for CPU-intensive work such as image processing or encryption.
      \item The asynchronous model achieves true parallelism by running tasks on multiple cores simultaneously.
      \item Threads generally incur higher overhead than async (creation, context switching, memory).
      \item The async model uses cooperative multitasking, where tasks voluntarily yield control to an event loop.
    \end{enumerate}

    \item \pts Select all that are TRUE:
    \begin{enumerate}
      \item At 60 FPS, the UI thread has roughly 16.67ms to produce each frame.
      \item Blocking the UI thread beyond the frame budget causes jank (dropped frames, stuttery animation).
      \item On Android, an unresponsive app can trigger an ANR (``Application Not Responding'') dialog after about 5 seconds.
      \item In Flutter, the thread responsible for UI is referred to as the Main/UI Isolate.
      \item CPU-intensive \texttt{async}/\texttt{await} code naturally yields back to the event loop, so it never blocks the thread.
    \end{enumerate}

    \item \pts Select all that are TRUE:
    \begin{enumerate}
      \item As a rule of thumb, an operation that takes more than \textasciitilde100ms of pure CPU time should be moved to an isolate.
      \item \texttt{compute()} creates an isolate, sends your data, receives the result, and then cleans up the isolate.
      \item You should generally not run more isolates than the device has CPU cores.
      \item Isolates can directly access and mutate the UI.
      \item Creating an isolate costs roughly 2--10ms.
    \end{enumerate}

    \item \pts Select all that are TRUE:
    \begin{enumerate}
      \item \texttt{setState()} tends to mix logic and UI together and does not scale well for complex apps.
      \item Lifting state up too far can lead to ``prop drilling'' --- passing data through many layers.
      \item Ephemeral (local) state lives within a single widget (e.g., a toggle switch or counter).
      \item \texttt{setState()} is the recommended approach for sharing the same state across many widgets and screens.
      \item App (global) state is shared across multiple widgets or screens (e.g., the logged-in user or theme mode).
    \end{enumerate}

    \item \pts Select all that are TRUE:
    \begin{enumerate}
      \item A Cubit emits new states directly, without separate events.
      \item A BLoC works with both events and states, offering more structure than a Cubit.
      \item Without Equatable, Dart's default equality compares object references rather than values.
      \item Equatable causes unnecessary rebuilds because it compares states by reference.
      \item BLoC stands for Business Logic Component and separates UI from business logic.
    \end{enumerate}

    \item \pts Select all that are TRUE:
    \begin{enumerate}
      \item Serverless (FaaS) bills per invocation and resource time, with no charge while idle.
      \item Cold starts add latency in serverless; mitigations include keeping minimum instances / provisioned concurrency.
      \item A VPS gives you full root control, but you are responsible for patching, firewalling, and key rotation.
      \item With a VPS you only pay for actual request processing and nothing while the machine is idle.
      \item Serverless typically wins for steady, constant high-traffic loads, while a VPS is better for spiky/low traffic.
    \end{enumerate}

    \item \pts Select all that are TRUE:
    \begin{enumerate}
      \item Idempotent HTTP methods (GET, HEAD, PUT, DELETE, OPTIONS) are defined to be safe to repeat.
      \item The \texttt{Retry-After} header tells the client how long to wait (e.g., after a 429 or 503).
      \item Exponential backoff combined with jitter helps avoid thundering herds and synchronized retries.
      \item Hedging is safe to apply to non-idempotent writes without any additional safeguards.
      \item ``Full jitter'' waits for a random value between 0 and the current exponential backoff delay.
    \end{enumerate}

  \end{enumerate}

  \newpage

  \nameheader
  \examtitle

%==================================================================
  \subsection*{Variant 2}
%==================================================================

  \answergridtwo

  \instructions

  \begin{enumerate}

    \item \pts A frame budget is $1000\,\text{ms} \div \text{target FPS}$. Assume the UI thread does all rendering work. Select all that are TRUE:
    \begin{enumerate}
      \item At 90 FPS the per-frame budget is about 11.11 ms ($1000 \div 90$).
      \item At 24 FPS the per-frame budget is about 41.67 ms ($1000 \div 24$).
      \item A synchronous task that runs for 250 ms on the UI thread blocks roughly 15 frames at 60 FPS ($250 \div 16.67 \approx 15$).
      \item At 120 FPS the per-frame budget is about 16.67 ms.
      \item A task that takes 8 ms at 60 FPS uses just under half of the frame budget ($8 \div 16.67 \approx 48\%$).
    \end{enumerate}

    \item \pts A client retries with exponential backoff: the base delay is \textbf{0.5 s} and it \textbf{doubles} every attempt (0.5, 1, 2, 4, \dots), with the delay for the \emph{n}-th retry computed as $0.5 \times 2^{\,n-1}$ seconds. Select all that are TRUE:
    \begin{enumerate}
      \item The uncapped backoff delay before the 5th retry is 8 s.
      \item With a 10 s maximum-delay cap, the 6th retry's delay is capped at 10 s instead of its uncapped value of 16 s.
      \item Using \textbf{full jitter} on a computed backoff delay of 4 s, the actual wait is a random value in the range [0 s, 4 s].
      \item Using \textbf{equal jitter} on a computed backoff delay of 4 s, the actual wait is a random value in the range [1 s, 2 s].
      \item Summing the uncapped delays of the first five retries (0.5 + 1 + 2 + 4 + 8) gives 15.5 s of total waiting.
    \end{enumerate}

    \item \pts Select all that are TRUE:
    \begin{enumerate}
      \item A web app is accessed using a URL and requires no installation.
      \item Web apps are cross-platform ``by default'' and usually require an active internet connection.
      \item Native iOS apps are built with Swift or Objective-C in Xcode.
      \item Cross-platform apps generally have worse performance than PWAs.
      \item A PWA is capable of working offline.
    \end{enumerate}

    \item \pts Select all that are TRUE:
    \begin{enumerate}
      \item An interpreted language executes code line-by-line via an interpreter at runtime.
      \item Compiled languages catch errors earlier, during compilation, preventing some runtime failures.
      \item Because interpreted source is often distributed alongside the interpreter, proprietary algorithms can be easier to reverse engineer.
      \item Compiled languages are better suited to rapid prototyping and scripting than interpreted ones.
      \item The performance gap matters most in computationally intensive applications (e.g., scientific computing, real-time systems).
    \end{enumerate}

    \item \pts Select all that are TRUE:
    \begin{enumerate}
      \item Dart is a client-optimized language developed by Google.
      \item Dart was first released in 2011.
      \item Dart is an object-oriented language.
      \item Dart cannot perform hot reload; it requires a full restart for every change.
      \item Dart compiles only via AOT and has no JIT mode.
    \end{enumerate}

    \item \pts Select all that are TRUE:
    \begin{enumerate}
      \item The young generation is collected independently and much more frequently than the old generation.
      \item For old-generation collections, the Dart VM uses concurrent marking that runs alongside the application (mutator) to reduce pause times.
      \item Incremental sweeping spreads the sweep phase across many small time slices to reduce long stop-the-world pauses.
      \item The young generation is collected less frequently than the old generation.
      \item A remembered set exists to track references from old-generation objects that point to young-generation objects.
    \end{enumerate}

    \item \pts A write barrier adds information to the remembered set, and the Dart VM has particular execution modes. Select all that are TRUE:
    \begin{enumerate}
      \item A write barrier executes whenever the application modifies an object reference.
      \item Both conditions---target object in the old generation AND new reference into the young generation---must hold for the barrier to record the reference.
      \item In JIT mode, Dart source code is initially parsed and executed through an interpreter.
      \item AOT-compiled output eliminates the VM overhead and produces standalone executables that start faster and use less memory.
      \item The Dart VM uses AOT mode during development and JIT mode in production.
    \end{enumerate}

    \item \pts Select all that are TRUE:
    \begin{enumerate}
      \item An isolate is Dart's approach to achieving true parallelism.
      \item Compared to Java/C++ threads (shared heap), isolates use separate heaps.
      \item Communication between isolates is done by direct memory access, just like threads.
      \item Synchronization primitives such as locks/mutexes are not needed between isolates.
      \item Large data can be moved between isolates using TransferableTypedData rather than a costly copy.
    \end{enumerate}

    \item \pts Select all that are TRUE:
    \begin{enumerate}
      \item In the threading model, code can run in parallel across multiple cores.
      \item The async model has lower overhead than threads because it avoids context switching between threads.
      \item Thread-based code is generally lower in complexity than async because there are no race conditions to manage.
      \item The async/await pattern looks synchronous while actually executing asynchronously.
      \item In the async model there are no shared-state issues because everything runs on the same thread.
    \end{enumerate}

    \item \pts Select all that are TRUE:
    \begin{enumerate}
      \item On iOS, the UI-handling thread is referred to as the Main Thread or Main Queue.
      \item Responsibilities of the UI thread include rendering, event handling, view lifecycle, and animations.
      \item On iOS, a watchdog can kill the app if it remains unresponsive.
      \item The golden rule is that heavy work should always be placed directly on the UI thread to keep it busy.
      \item Dropped frames caused by exceeding the frame budget are commonly called ``jank.''
    \end{enumerate}

    \item \pts Select all that are TRUE:
    \begin{enumerate}
      \item If an operation primarily waits for I/O (network, disk), async/await is preferred over an isolate.
      \item \texttt{compute()} also handles errors, in addition to creating the isolate, sending data, receiving the result, and cleaning up.
      \item Sending a small message ($<$ 1 MB) between isolates is a fast copy.
      \item Isolates have effectively zero communication overhead because there is no serialization cost.
      \item On a 4-core device you should generally not run more than 4 isolates.
    \end{enumerate}

    \item \pts Select all that are TRUE:
    \begin{enumerate}
      \item In Flutter, widgets are immutable once created; when state changes, Flutter rebuilds widgets with the new data.
      \item Immutability aids predictability because the UI does not mutate unexpectedly.
      \item Flutter uses keys to compare old and new widgets, keeping elements that did not change and rebuilding those that did.
      \item A recommended best practice is to keep UI (widgets) and business logic tightly coupled in the same place.
      \item A reasonable progression is to start simple with \texttt{setState} and move to Riverpod/BLoC as complexity grows.
    \end{enumerate}

    \item \pts Select all that are TRUE:
    \begin{enumerate}
      \item BLoC uses a stream-based approach with events and states via Stream and Sink.
      \item In the BLoC architecture, the data layer handles API calls and databases.
      \item Equatable lets Flutter compare states by value, avoiding unnecessary rebuilds.
      \item Without Equatable, two \texttt{CounterState(1)} objects with the same count are treated as equal by Dart's default equality.
      \item A Cubit is a simpler version of BLoC that emits new states directly without separate events.
    \end{enumerate}

    \item \pts Select all that are TRUE:
    \begin{enumerate}
      \item Serverless platforms handle provisioning, scaling, and patching for you.
      \item A benefit of serverless is automatic scaling, from zero up to peaks, with a small operations surface.
      \item A VPS such as Hetzner Cloud can use Floating IPs that move between nodes to support high availability.
      \item Serverless functions have a stateless runtime and execution time limits.
      \item Serverless gives you full root control of the kernel and installed agents.
    \end{enumerate}

    \item \pts Select all that are TRUE:
    \begin{enumerate}
      \item POST and PATCH should be retried with caution unless you add idempotency keys.
      \item Hedged requests send a second copy of the same request after a short delay and use the first successful response, canceling the rest.
      \item gRPC supports a hedging policy natively.
      \item Decorrelated jitter randomizes delays around a moving average to adapt faster.
      \item Hedging should be applied freely to non-idempotent writes since it always reduces latency safely.
    \end{enumerate}

  \end{enumerate}

\end{document}
